\section{Software de mezclado de vídeo en directo}

Antes de abordar el diseño de Voctomix 2.0 se realizó un análisis de las herramientas
de mezcla de vídeo en directo disponibles, tanto de código abierto como comerciales,
con el fin de identificar las carencias que justifican las aportaciones de este trabajo.

\textit{dvswitch} fue durante años la herramienta estándar de C3VOC. Su principal
limitación es la restricción al formato DV (720×576 SD), que la hace inviable para
producciones en alta definición. Además, su arquitectura monolítica no permite separar
el motor de procesamiento de la interfaz de operador.

\textit{OBS Studio} es la herramienta de streaming de código abierto más popular en
la actualidad. Ofrece overlays, transiciones y una interfaz moderna, pero su
arquitectura está orientada a un único operador en un único equipo: no existe un modo
cliente-servidor nativo que permita separar el procesamiento de la interfaz.

\textit{vMix} es una solución comercial profesional para Windows con soporte para NDI,
multiview y producción 4K. Su coste de licencia y su limitación a Windows la excluyen
de entornos de producción \textit{open-source} sobre Linux.

Voctomix 2.0 se posiciona en el espacio que ninguna de estas herramientas cubre de
forma satisfactoria: arquitectura cliente-servidor nativa, \textit{scripting} TCP,
overlays dinámicos, soporte multi-audio y despliegue reproducible mediante Docker,
todo ello bajo licencia GPL.

\subsection{Voctomix original}

Voctomix fue concebido por el equipo C3VOC como respuesta a las limitaciones de
\textit{dvswitch}, el mezclador empleado hasta entonces en el Chaos Communication Congress.
\textit{dvswitch}, aunque funcional para resoluciones estándar de definición estándar (SD),
carecía de soporte nativo para alta definición (1920×1080) y no permitía separar el motor
de mezcla de la interfaz de operador, lo que dificultaba su uso en producciones distribuidas
donde el realizador trabaja en un equipo diferente al servidor de procesamiento.

La arquitectura de Voctomix se articula en torno a dos procesos independientes que
se comunican mediante sockets TCP: \texttt{voctocore}, responsable de todo el
procesamiento multimedia, y \texttt{voctogui}, la interfaz gráfica de control.
El núcleo acepta fuentes de vídeo en formato UYVY crudo y audio PCM S16LE
encapsulados en contenedores Matroska sobre TCP, los mezcla mediante el compositor
de GStreamer y expone el resultado en distintos puertos de red para su consumo por
parte de la GUI, grabadores o encoders de streaming.

El protocolo de control (puerto 9999) define un conjunto de comandos de texto plano
(\texttt{set\_video\_a}, \texttt{set\_composite\_mode}, \texttt{set\_stream\_blank}, etc.)
que cualquier cliente TCP puede enviar, lo que facilita la automatización y la integración
con hardware externo como controladores MIDI, paneles de control físicos o scripts de
producción automatizada.

% --- Tabla sugerida ---
% \begin{table}[H]
%   \centering
%   \caption{Comparativa de herramientas de mezcla de vídeo en directo.}
%   \label{tab:comparativa_herramientas}
%   \begin{tabular}{|l|c|c|c|c|c|c|}
%     \hline
%     \textbf{Herramienta} & \textbf{Licencia} & \textbf{Arch.} & \textbf{Multi-PC} & \textbf{Overlays} & \textbf{Docker} & \textbf{Scripting} \\
%     \hline
%     dvswitch       & GPL        & Monolítica       & No      & No  & No  & No \\
%     OBS Studio     & GPL        & Monolítica       & Parcial & Sí  & No  & Sí (Lua) \\
%     vMix           & Comercial  & Monolítica       & No      & Sí  & No  & Sí \\
%     \textbf{Voctomix 2.0} & \textbf{GPL} & \textbf{C/S} & \textbf{Sí} & \textbf{Sí} & \textbf{Sí} & \textbf{Sí (TCP)} \\
%     \hline
%   \end{tabular}
% \end{table}
